\subsection{Breakpoints: Inspecting a Suspended Lexical Environment}
\label{sec:break}

Retry repaired a precondition and resent a request; here the application
deliberately waits while we inspect its current bindings.
\filepath{examples/misc/break.abclr2} implements an interactive debugger
in ABCL/R2 itself. It does not invoke the Common Lisp debugger or install
an operating-system interrupt handler. An application calls a macro,
waits for a reply, and lets another object inspect its lexical environment
before execution continues. This makes reflection useful for a concrete
task rather than merely printing an object's representation.

\subsubsection{Reproduce the Original Session}

Load the debugger before defining code that uses its macro. At the repository
root, enter the following. The definition of \code{foo} comes from the
original file's commented session, not from a replacement example.

\begin{codeblock}
\begin{lstlisting}
(inc "examples/misc/break" :use-source t)
[object foo
 (state [x := 0])
 (script
  (=> y
      (temporary [z := (+ x y)])
      (abclr2-break)
      !z))]
[foo <== 123]
\end{lstlisting}
\end{codeblock}

The handler binds \code{y} to 123 and computes the temporary \code{z}
as \code{x + y}, initially 123. The breakpoint now prints an object
identifier and an \code{Inspect(...)} prompt. Object numbers and generated
reply-variable names vary between sessions; they are not expected results.
At this prompt you are talking to the environment inspector, not the
ordinary \code{ABCL/R2>} reader.

Use the following sequence. For \code{:set}, the next two inputs answer
the \code{Name =} and \code{Value =} prompts.

\begin{codeblock}
\begin{lstlisting}
:l             ; lists Z in the temporary frame
z              ; prints Z=123
:u
:l             ; lists Y and a generated reply binding
y              ; prints Y=123
:u
:l             ; lists X in the application-state frame
:set
x
10
:d
:d
:set
z
999
:r
0
\end{lstlisting}
\end{codeblock}

The first \code{:u} reaches the handler's pattern bindings, and the second
reaches the object's state. The two \code{:d} commands retrace those steps.
The resumed handler sends 999, because you changed its current \code{z}.
The value 0 supplied after \code{:r} is the result of
\code{(abclr2-break)}, which this handler ignores before executing
\code{!z}.

Now send \code{[foo <== 123]} again. At the new inspector prompt, enter
\code{:r}, then \code{t}. The handler returns 133: \code{x} remained 10,
but a new invocation computed a new temporary \code{z = 10 + 123}.
This is the original demonstration's distinction between persistent state
and an invocation-local binding. The inspector changes existing bindings;
it does not automatically recompute an already calculated \code{z}
when you change \code{x}.

\subsubsection{The Inspector's Small Command Language}

\begin{longtable}{L{0.25\textwidth}L{0.65\textwidth}}
\textbf{Input} & \textbf{Effect}\\\hline
\endhead
\code{:l}, \code{:list} & List the names in the selected local frame.\\
\code{:a}, \code{:alist} & Print that frame's association list of names and values.\\
A variable name & Look up its value, searching outward if necessary;
print an unbound notice if lookup fails.\\
\code{:u}, \code{:up} & Select the outer lexical environment, if one exists,
and remember the current frame.\\
\code{:d}, \code{:down} & Return to the most recently remembered frame,
if there is one.\\
\code{:s}, \code{:set}, \code{:=} & Read a name and a value, then update
or create that binding in the selected local frame.\\
\code{:h}, \code{:help}, \code{?} & Print the built-in command help.\\
\code{:r} & Read a return value, reply to the breakpoint caller, and finish
the inspector handler.\\
\end{longtable}

These commands are read by Lisp's \code{read}, but are interpreted by
the inspector's \code{case}, not evaluated as arbitrary Lisp forms.
Enter \code{x}, not \code{'x}, when asked for a name. Entering
\code{(+ 1 2)} as a value supplies a list; it does not calculate 3.
In particular, \code{:r} is a command followed by a separate value input.
It is missing from the original printed help, although the loop's
termination test implements it.

Local listing and outward lookup are different operations. You can look up
\code{x} while positioned at the temporary frame even though \code{:l}
there lists only \code{z}. However, \code{:set} uses \code{:set-local},
not the outward-searching \code{:set} environment operation. Setting
\code{x} in the temporary frame would introduce a local binding there,
not update the object's existing state. Navigate to the state frame first
when you intend the persistent change shown above.

\subsubsection{How the Macro Obtains the Environment}

Read the original monitor:

\wholelisting[firstline=7,lastline=16]{examples/misc/break.abclr2}

The macro at lines 65--66 expands into a now-type request to this monitor.
It creates an anonymous, empty-script \emph{dummy object} at the call site
and includes that object, an optional format string, and a list of format
arguments in the message. The macro uses backquote, comma, and
comma-at as introduced in the Lisp reading guide: it builds code while
inserting the caller's forms.

Why create an object that handles no messages? Its state environment
retains the lexical environment in which it was created.
\code{[meta dummy-object]} obtains its metaobject;
the \code{:state} query obtains its state environment; and
\code{:lexenv} follows that environment's outer link. That link is the
application environment at the breakpoint. It is a live environment object,
not a copied map or a snapshot of a machine stack.

The monitor receives the breakpoint request's reply destination as
\code{R} and its application sender as \code{sender}. It prints that sender,
then sends \code{[environment nil]} to \code{environment-inspector},
forwarding \code{@ R}. The initial \code{nil} is an empty navigation
history. The monitor does not manufacture the final return value or wait
for a second result: the inspector replies directly to the waiting caller.
If a format string is present, \code{~?} applies it to the supplied
argument list to print a contextual message.

\subsubsection{Navigation, Mutation, and Resume}

The inspector explicitly selects \code{non-reifying-meta}. Its handler is
compiled for the non-reifying execution path rather than interpreted by a
custom reflective evaluator. It can still send messages to the environment
objects it was given. Non-reifying does not mean that the supplied
environment has become inaccessible.

\wholelisting[firstline=21,lastline=29]{examples/misc/break.abclr2}

The loop begins with the command \code{:skip}, whose case does nothing.
Its step expression prints the prompt and reads the next command.
The termination test recognizes \code{:r}; it reads one more datum and
sends it to \code{R}. Until this happens, the application's now-type
breakpoint request has not received its reply. The suspended application
handler resumes after that reply, not after a local-listing or edit command.

\wholelisting[firstline=31,lastline=49]{examples/misc/break.abclr2}

For upward navigation, \code{env-record} is a stack of previously selected
environments, represented as a Lisp list. The bracket expression
\code{[environment . env-record]} constructs a new head/tail pair.
Downward navigation retrieves its \code{car} and removes it with
\code{cdr}. It therefore retraces your own upward moves; it does not
search for every inner scope in the program. An upward move with no outer
environment or a downward move with empty history does nothing.

The setting branch uses \code{let*} so the name is read before the value,
then sends the local-binding update asynchronously. The original session
and isolated checks show those updates visible when the handler resumes.
Do not extend that observation to arbitrary debugger clients or scheduler
policies without checking their message ordering.

\paragraph{Evidence and scope.}
Run \code{make -C tutorial break-check}. Six checks replay the original
definition and editing session from unchanged source copies, verify
999 followed by 133, inspect persistent \code{x}, and confirm that
a supplementary handler receives \code{(+ 1 2)} as raw breakpoint-return
data. Transcript assertions also check local listings, variable values,
unbound-name handling, help, and an optional formatted message.
They do not certify every command alias or concurrent breakpoint arrival.
The single inspector reads shared interactive input; this example is not
a multi-user debugger, a rollback facility, or a whole-system pause.

\paragraph{Check your understanding, with answers.}
Why does the second call not return 999?
The first call's \code{z} was temporary; the second creates a fresh binding.
Why can its result still depend on the first session?
The state binding \code{x} was changed in the persistent state frame.
Does entering \code{:r} and then 0 force the application to return 0?
No. It returns 0 from the breakpoint expression; the handler still decides
what application reply to send afterwards.
